\section{Discussion}
\label{sec:discussion}

\paragraph{What decision models offer as conflict classifiers.} Conflict typing sits on the critical path of every
search-augmented answer, so its latency and cost add to every request. Most hosted decision models classify a query
with about ten retrieved pages in under half a second, for at most \$\dmcostmax per thousand queries, with the accuracy of the LLM used in the original study. When the accuracy of a
current LLM is required, a decision model's confidence is usually good enough to route only its uncertain queries to
the LLM, which recovers most of the gap at a quarter to a third of the LLM's cost. End to end, the cheaper types cost
no measurable answer quality (\S\ref{sec:endtoend}). 

\paragraph{Choosing among decision models.} The class is not uniform. Accuracy spans \dmaccmin--\dmaccmax\%, and the
model at the bottom is significantly worse than the paper's LLM. Two properties that matter in deployment vary
independently of accuracy: calibration (ECE \eceminimum--\ecemaximum) and how well confidence ranks errors, which
decides whether a cascade helps at all. One of the most accurate models has confidence that is no better than random for
escalation. A model for this task should therefore be chosen on all three, measured on the deployment's own data,
not on accuracy alone.

\paragraph{How to ask.} The direct choice question was the reliable format for every model. The yes/no
decomposition looked natural but failed for all of them; the repaired tree encodes choices fitted to one model on a
small development split and did not help any other. For a new decision model we would start from the flat
question.

\paragraph{What the diagnostics change.} Dates turned out not to matter: outdated sources are recognised
from their content. Sensitivity to source order is real but differs so much between models that it has to be
measured for each one. A single false source is the subtle case, visible in the probabilities but rarely in the
top label. With
\nmisitest natural \emph{Misinformation} test queries, that type is the one a deployment would most need to
calibrate on its own data.

\paragraph{Limitations.}
(i) \conflicts is small: with \ntest test queries, differences under about five points are not significant, and the
cascades can be shown close to, but not formally equivalent to, the LLM. (ii) The label distribution reflects how the
queries were selected, not real search traffic, and the dataset has only five natural \emph{Misinformation}
queries; our planted variant tests one kind of misinformation (an entity or value substituted in a single source).
(iii) The census covers the decision models one gateway served on one day. Hosted models change without notice:
several identifiers are unversioned, two models were free during our runs, and our rerun of Gemini~2.5 Flash scores
below the published figure. Providers
do not disclose model sizes or training data, so we cannot explain the differences between models. (iv) The end-to-end test uses one generator and LLM judges, validated on \nvalid answers, that do not see the
search results; as in the original study, the primary judge is the generator's own model, which may favour its
answers \citep{panickssery2024llm}. The size of the gain from predicted types depends on the judge, and we do not
measure factual grounding, the third metric of \citet{cattan2025dragged}. (v) Not every system is deterministic: identical requests changed
up to \rerunflipmax\% of a decision model's predictions, and \diaggmrerunflip\% of Gemini~3.8 Flash's.
